\documentclass[11pt]{article}
\usepackage[margin=1in]{geometry}
\usepackage{booktabs,amsmath,hyperref,xcolor,graphicx}
\usepackage[numbers]{natbib}
\newcommand{\todo}[1]{\textcolor{red}{[TODO: #1]}}

\title{Androids Dream of the Sheep They Saw That Day:\\
Test-Time Memory and Dream Consolidation in a Small Language Model}
\author{Gianni Sanrochman\\\texttt{gsanrochman@proton.me}}
\date{November 2026 \quad (draft)}

\begin{document}
\maketitle

\begin{abstract}
We build and measure a complete learn-by-day, consolidate-by-night loop in a
134M-parameter language model. During the day the model reads a text stream
with its weights frozen; new facts live only in a fast tier, a cache of the
previous window's attention keys and values. At night the model is cued with
answer-free prompts about the day and completes them from that cache. These
\emph{dreams}, filtered by self-consistency voting, are the only training
data for a set of sparse lifetime tables; a gate on held-out loss and on
recall of everything committed before can roll the write back. No ground-truth
answer enters consolidation. Three results stand out. (1) A retrieval
curriculum that keeps every answer out of the attention window, the standard
guard against leakage, trained the model out of in-context copying (31.8\% to
3\% first-token accuracy with the fact in context), silently confounding four
rounds of memory-architecture experiments. (2) Every newly added memory module
we tried (a TTT/Titans-style fast-weight matrix, a Nested-Learning continuum
memory, per-token write strength, an uncompressed slot memory) failed at
\emph{addressing} rather than storage; reusing the backbone's own attention
through a Transformer-XL cache raised one-window-back recall from 4.3\% to
94.1\%. (3) Voted dreams consolidate at 81\% of the ground-truth control and
preserve old knowledge under later learning. (4) The first full loop cleared
its pre-registered bars only at the margin, because the cache could hold a fact
without the model being able to answer a cue from it (27\% at realistic
density). Short training on the cue format with held-out subjects and
per-sample random answers raised cache dreams to 100\%; tuning only the
attention query and key projections, against real post-sleep tables, does so
within 0.013 of the untuned model's val loss, and the loop then learns 57\% of a day's
facts (0\% without sleep) and keeps half of them after two more days. All
thresholds were set before each run, and every probe is reported against the
model's own in-context ceiling.
\end{abstract}

\section{Introduction}
A deployed language model that learns something new must either keep it in
context, which is bounded, or change its weights, which interferes with what it
already knows. Complementary learning systems theory
\citep{mcclelland1995cls} describes how brains split the problem: the
hippocampus records episodes quickly, and sleep replays them into slow
cortical memory. Recent architectures import the fast half of that story as
test-time learning \citep{sun2024ttt,behrouz2025titans,behrouz2025nested}.
The slow half, and the loop between them, is less explored in language models.

We build the whole loop at small scale and measure each stage separately before
measuring it end to end. Our model, \emph{accretion}, has three tiers:
\begin{itemize}
\item a \textbf{fast tier} that holds what the model read today without any
gradient step;
\item \textbf{lifetime tables}, sparse hashed embeddings injected into the
residual stream and written only during sleep;
\item a \textbf{sleep job} that writes the tables from the model's own
\emph{dreams}, rehearses older knowledge by re-dreaming it, and is guarded by
an external gate with rollback.
\end{itemize}

The paper's contributions are as much about measurement as about the system.
Two of our intermediate conclusions were wrong for the same reason: a pass bar
was set without first measuring what the unassisted backbone could do with the
answer in plain view (Section~\ref{sec:calibration}). We report those errors,
because the failure mode they expose is general.

\section{System}
\subsection{Backbone}
A 24-block decoder, $d=640$, GPT-2 tokenizer, trained on FineWeb-10B
\citep{penedo2024fineweb}. Eighteen blocks are Gated DeltaNet
\citep{yang2025gdn}; every fourth block is grouped-query attention with QK-norm
and no positional encoding. An ablation queue at 2B tokens selected this hybrid
over dense attention (val 3.259 vs 3.286) and rejected multi-token prediction
and a content-blind Engram-style table. Optimizer: Muon for matrices, AdamW
elsewhere, warmup-stable-decay schedule.

\subsection{Fast tier}
We tested five fast tiers, each trained jointly with the backbone on a
retrieval curriculum (Section~\ref{sec:curriculum}):
a TTT-Linear fast-weight matrix with learned initial state and momentum
\citep{sun2024ttt,behrouz2025titans}, gated into block 12 (MAG);
a three-level continuum memory system at update periods 32/128/512 tokens
\citep{behrouz2025nested};
backpropagation through two consecutive windows plus an InfoNCE readback loss;
per-token write strength; and an uncompressed slot memory that stores every
token's key and value and reads them by multi-head attention.
The tier that works is the simplest: a Transformer-XL cache
\citep{dai2019transformerxl} (no new parameters beyond one bias per key-value head). Each attention block also attends to the previous
512-token window's keys and values, placed before the current window; a
learned per-head bias on cached logits marks them as past
(Section~\ref{sec:results-fast}).

\subsection{Lifetime tables and sleep}
Lifetime memory is a pair of hashed $n$-gram tables injected after blocks 2 and
12. A sleep cycle trains only the tables (AdamW, 12k steps) on restructured
documents about the day's facts. Two safeguards come from earlier failures.
\emph{Replay}: 30\% of each batch rehearses previously committed facts.
\emph{Gate}: a write is committed only if held-out val loss stays within
$+0.05$ of the original baseline (not of the previous cycle, which let small
drifts compound to $+0.09$) and previously committed facts keep at least half
of their recall; otherwise the pre-sleep snapshot is restored. The model never
certifies its own edit.

\subsection{Dreaming}
\label{sec:dreaming}
In the ground-truth version of sleep the training documents contain the true
answers. In the dreaming version they never do. For each fact the registry holds
only a cue index (subject and relation). Each of seven restructured phrasings
is cut just before the answer, producing a cue such as
\texttt{Q: Who founded Quonsirn?\textbackslash nA:}; the object is replaced by
a sentinel before templating, so it cannot leak (tested). The model completes
each cue greedily. The day's facts are dreamed with the day's memory available;
older facts are re-dreamed from the current weights alone. With
\emph{consensus}, the seven dreams vote on the answer and every cue is
completed with the winner, which also strips invented detail. Wrong dreams are
consolidated anyway and counted as false memories.

\section{Benchmarks and calibration}
\label{sec:calibration}
\subsection{Cross-window recall}
\label{sec:curriculum}
A quarter of training lanes carry a synthetic stream: novel facts with random
two-syllable names (``the mayor of Blafro is Kresut''), filler sentences, and
probes asking for a fact one or more windows after its last mention. An
anti-leak invariant guarantees that no probe shares a window with any mention
of its fact. The metric is first-token accuracy on the answer, with memory
carried across windows minus the same checkpoint with memory reset each window,
bucketed by the probe's age in windows.

\subsection{The ceiling we did not measure}
We later measured the same probes with the statement placed back in context
and memory off (Table~\ref{tab:ceiling}). Before curriculum training the
backbone copies a novel name from a nearby sentence 31.8\% of the time, falling
to chance by 256 tokens (2.7\%, against 2.5\% with no fact in context). After curriculum training it does not copy at all:
3.0\%, equal to the no-fact floor. Because the anti-leak rule kept every answer
out of attention, and the untrained memory supplied nothing, the best policy
the model could learn was to guess a plausible name. Every memory arm trained
on that curriculum was feeding an answer path the curriculum had switched off.
Re-asking 15\% of facts inside their own window (``near'' probes, logged
separately and excluded from scoring) restored the circuit to 95\% and extended
it to 89\% at 450 tokens.

\begin{table}[t]\centering\small
\caption{First-token accuracy, fact in context, memory off, by distance in tokens
between statement and question (403 probes).}\label{tab:ceiling}
\begin{tabular}{lrrrrrrrr}\toprule
checkpoint & 0 & 16 & 32 & 64 & 128 & 256 & $\sim$450 & no fact\\\midrule
base (no curriculum) & 31.8 & 26.1 & 22.3 & 12.9 & 7.9 & 2.7 & 2.7 & 2.5\\
old curriculum & 3.0 & & & & & & 4.2 & 5.0\\
with near probes & 95.0 & 95.5 & 95.5 & 95.5 & 95.0 & 94.8 & 89.1 & 6.7\\\bottomrule
\end{tabular}\end{table}

\subsection{Lifetime benchmark}
Three disjoint corpora of 60 facts each are consolidated in sequence; after
every cycle every earlier corpus is re-probed with a light paraphrase whose
in-context ceiling is 100\%. Answers come from a list of 10 names, so random
guessing scores 10\%. Gate-off runs are mandatory: a gate that rolls everything
back forgets nothing and learns nothing.

\section{Results}
\subsection{The fast tier fails at addressing}
\label{sec:results-fast}
Table~\ref{tab:fast} lists every fast-tier round. Rounds 0 to 2 are confounded
by the curriculum. Round 3 is the first fair test of the fast-weight matrix:
+3.5pp. A diagnostic that ranks the stated fact's tokens among all keys of the
previous window shows why: the query at the question finds the fact only
diffusely (median best rank 18 of 512) and the read returns a blend.

Round 4 removed compression entirely. The slot memory's gate opened fully (1.00)
and its output projection's norm grew about 300-fold, so the model used
the branch heavily, but not for facts: attention mass on the fact's tokens was
3.9\% against 2.7\% chance. A freshly initialized memory must learn retrieval
from sparse probe gradients and instead learned to inject generic context. The
backbone's own attention already solves the task at 89 to 95\% in context.
Round 5 reuses it, and recall one window back rises to 94.1\% against 4.3\%
without the cache. Ages 2 and beyond stay at zero, as expected for a one-window
cache.

The cache has a cost. With the cache carried, FineWeb val is 3.498 against
3.342 for the round-3 model; the gap persists under every evaluation protocol.
Our hypothesis was that NoPE attention infers order from the causal mask, which
an always-visible cache breaks. Round 6 tested it with a learned per-head bias
on cached logits, plus a denser curriculum (50\% fact lines against 35\%) with
facts wrapped in document-style framings. Recall one window back held (+88.4),
but val did not move (3.496) and the learned bias stayed near zero, so the
hypothesis is not supported. The gap remains open; we leave it as a
limitation.

\begin{table}[t]\centering\small
\caption{Fast-tier rounds. Memory gain = first-token accuracy on facts stated one
window (512 tokens) earlier, memory carried minus reset, same checkpoint
(423 probes; SE of a difference $\approx$2.3pp). Pre-registered pass: $\geq$+10pp.}
\label{tab:fast}
\begin{tabular}{llrl}\toprule
round & fast tier & gain & note\\\midrule
0 & TTT-Linear matrix & +1.4 & \\
1 & 3-level CMS & $-$0.2 & slow levels pruned; val 3.94\\
1 & matrix + 2-window BPTT + readback & $-$0.9 & readback 1.01 vs chance 4.16\\
2 & + per-token write strength & $-$1.7 & \\\midrule
\multicolumn{4}{l}{\emph{curriculum fixed: in-context copying restored}}\\
3 & per-token matrix & +3.5 & read rank 18/512\\
4 & slot memory (no compression) & +1.2 & gate 1.00; fact mass 3.9\% vs 2.7\%\\
5 & Transformer-XL cache & \textbf{+89.8} & 94.1\% vs 4.3\%; val +0.16\\
6 & XL + cache bias + denser, wrapped curriculum & +88.4 & val 3.496; bias $\approx$0\\
7c & round 5 + cue tuning (q/k only) & +88.2 & val 3.511\\\bottomrule
\end{tabular}\end{table}

\subsection{Sleep, replay and interference}
On the earlier memory-native checkpoint, consolidation without replay erased
corpus A almost entirely after two later cycles (kept 12.5\% homogeneous, 0\%
on varied domains). With 30\% replay, corpus A kept 77.8\% and 76.5\%, and total
retained knowledge rose 2$\times$ and 6 to 8$\times$. A 2$\times$2 places the
interference in shared projection capacity rather than bucket collision: varied
corpora with 1.3\% bucket overlap still forgot completely, while homogeneous
corpora with 56\% overlap retained more; freezing the shared projections
delayed collapse by one cycle and crippled new learning. Replay works because
it keeps old facts in the gradient that shapes the shared projections.

A caution on the gate: the $+0.05$ drift budget sits at run-to-run noise. The
same ground-truth configuration committed its first cycle at $+0.045$ in one run
and rolled it back at $+0.051$ in another (the tables are randomly
initialized). We therefore report gate-off runs as the primary measurement.

\subsection{Dreaming}
Table~\ref{tab:dream} compares sleep on the true answers with sleep on the
model's own dreams. Raw dreams preserve old knowledge nearly as well as the
truth (corpus A keeps 75\%) but consolidate only 49\% as many facts and drift
past the budget. The dreams look like dreams: the core fact is usually right,
details are invented and inconsistent (``founded by Nimbleton in 1855'', then
``in 1844''), and some are false (``the capital of Lutvay is Lutvay'').
Consensus voting fixes both failures: 81\% of the control's retained facts
($\sim$79 against $\sim$98 of 180) and drift within budget (+0.043). Re-dreamed old facts are only as accurate as current
recall, so mis-remembered facts are rehearsed as faithfully as correct ones.

\begin{table}[t]\centering\small
\caption{Sleep on ground truth vs on dreams (memory-native checkpoint, day held
in context, gate off, 3$\times$60 facts, 30\% replay).}\label{tab:dream}
\begin{tabular}{lrrr}\toprule
& truth & raw dreams & voted dreams\\\midrule
facts retained of 180 & $\sim$98 & $\sim$48 & $\sim$79\\
corpus A kept after 2 more cycles & 81\% & 75\% & 58\%\\
val drift (budget +0.05) & +0.041 & +0.058 & +0.043\\
day dream accuracy & -- & 56--67\% & 77--100\%\\\bottomrule
\end{tabular}\end{table}

\subsection{The full loop}
On the round-5 model, 60 facts are packed whole into six 512-token windows with
filler and streamed through the model with the cache carried. Each fact is
dreamed from the cache state right after its window, with its text no longer in
context. Voted dreams train the tables; recall is probed on the following days
with the cache empty (Table~\ref{tab:loop}).

\begin{table}[t]\centering\small
\caption{Full loop on the round-5 model, gate off. Recall is with the fast tier
empty.}\label{tab:loop}
\begin{tabular}{lrrr}\toprule
& truth & dreams, day in context & dreams, day in cache\\\midrule
recall with no sleep & 0--3\% & 0--3\% & 0--3\%\\
corpus A after its sleep & 58.3\% & 43.3\% & 10.0\%\\
corpus A after 2 more days & 58.3\% & 18.3\% & 5.0\%\\
facts retained of 180 & $\sim$101 & $\sim$68 & $\sim$24\\
day dream accuracy & -- & 80--98\% & 15--35\%\\
val drift & +0.010 & +0.018 & +0.018\\\bottomrule
\end{tabular}\end{table}

The loop clears its pre-registered bars (learning $\geq$+10pp over no sleep;
corpus A keeps $\geq$50\%; drift $\leq$+0.05), but at the threshold: six facts
learned, three kept (60 facts per corpus, one seed). Every stage behaves as measured in isolation. Sleep turns
between half and two thirds of correct dreams into lasting recall in every
arm (10/15, 22/35, 20/29 from the cache; 43/80 with the day in context), so the bottleneck is the first
stage. Dream accuracy from the cache falls with fact density (voted: 70, 55,
47, 33, 27\% at 1, 2, 3, 5 and $\sim$10 facts per window), and even one fact per
window reaches only 70\% on framings the model never trained on.

Round 6 retrained the backbone on denser, document-wrapped facts to close that
gap. It made it worse: voted cache dreams fell to 52, 37 and 8\% at 1, 3 and
packed facts per window, and the loop failed three of its four bars (corpus A
+8.3pp, kept 40\%). Denser retrieval training teaches the model to answer
\emph{probes} from the cache; it does not teach it to answer \emph{cues}.

\subsection{Cue tuning}
\label{sec:cuetune}
We therefore trained the dreaming step directly. Each sample is two windows.
The first holds 1 to 12 fact documents in the format the day stream uses,
padded with FineWeb filler to 512 tokens; the second holds the night's cues for
those facts, each followed by its answer, with the first cue at position 0 so
that nothing of the first window is in context. The loss is cross-entropy on
the answer tokens only, plus the ordinary language-modeling loss on a FineWeb
batch at weight 1, and gradients flow through the carried cache. Two choices
keep this from being memorization. The subjects of the three evaluation corpora
are held out of training (241 of the 400 possible names remain), and every
answer is drawn fresh for every sample, so no subject is tied to any object in
the weights; the only thing that can be learned is to read the answer from the
cache. Leak controls on the evaluation corpora confirm it: voted dream accuracy
is 100\% with the right cache, 13 to 20\% with another day's cache, and 0 to
2\% with no memory (chance is 10\%).

Table~\ref{tab:cuetune} lists the variants. Tuning all weights (round 7, lr
$5\times10^{-5}$, 1600 steps) gives perfect cache dreams at every density and
the strongest loop, but it costs language modeling (val carried 3.673, against
a pre-registered ceiling of 3.52). Tuning only the query and key projections of
all attention layers (2.9M parameters) keeps val within budget (3.508) and keeps
the cache dreams, but corpus A keeps only 32\%: dreams of the current day decay
from 85 to 66\% across nights. The cause is that the cue read was tuned with the
lifetime tables empty, and after sleep they are not. Round 7c continues the
query/key tuning for 800 steps with the tables sampled from three real
post-sleep states (built from training-name corpora only). It passes all four
pre-registered bars: val 3.511, cache dreams 100\% at every density, corpus A
learned at 56.7\% against 0\% without sleep, kept at exactly 50.0\% after two
more days, drift +0.024.

\begin{table}[t]\centering\small
\caption{Full loop after cue tuning (day streamed through the cache, voted
dreams, gate off, 3$\times$60 facts, one seed). ``Truth'' trains the round-7
model's sleep on the true answers. Bars: val $\leq$3.52, A learned $\geq$+20pp
over no sleep, A kept $\geq$50\%, drift $\leq$+0.05.}\label{tab:cuetune}
\begin{tabular}{lrrrrr}\toprule
& r5 & r7 all & r7 truth & r7b q/k & r7c q/k + tables\\\midrule
val, cache carried & 3.498 & 3.673 & 3.673 & 3.508 & \textbf{3.511}\\
cache dreams, packed & 27\% & 100\% & -- & 100\% & 100\%\\
corpus A after its sleep & 10.0\% & 56.7\% & 58.3\% & 56.7\% & 56.7\%\\
corpus A after 2 more days & 5.0\% & 38.3\% & 55.0\% & 18.3\% & 28.3\%\\
A kept & 50\% & 67.6\% & 94.3\% & 32.3\% & \textbf{50.0\%}\\
facts retained of 180 & $\sim$24 & $\sim$79 & $\sim$100 & $\sim$46 & $\sim$56\\
day dreams, nights 1/2/3 & & 100/95/88 & -- & --/85/66 & 100/93/80\\
replay dreams, nights 2/3 & & 60/42 & -- & 28/23 & 47/24\\
val drift & +0.018 & +0.022 & +0.020 & +0.021 & +0.024\\\bottomrule
\end{tabular}\end{table}

The remaining loss is in rehearsal. Old facts are re-dreamed from the weights,
and by night 3 only 24\% of those dreams are right, so most of what is rehearsed
about corpus A is wrong. \todo{round 8: vote-agreement gate on replay dreams}

\section{Related work}
\textbf{Test-time memory.} TTT layers \citep{sun2024ttt} and Titans
\citep{behrouz2025titans} learn a fast-weight memory during the forward pass;
Nested Learning and Hope \citep{behrouz2025nested} add memories at several
update frequencies. Gated DeltaNet \citep{yang2025gdn} uses a delta-rule
state with per-token gates. Our negative results for these mechanisms are
specific to a 134M model, a sparse retrieval signal and cross-document state;
they say the bottleneck was learning to address, not capacity.
\textbf{Retrieval over cached states.} Transformer-XL \citep{dai2019transformerxl}
and Memorizing Transformers \citep{wu2022memorizing} attend to cached keys and
values; MemoryLLM and M+ \citep{wang2024memoryllm,wang2025mplus} maintain a
self-updated latent memory pool in an 8B model.
\textbf{Consolidation and replay.} Complementary learning systems
\citep{mcclelland1995cls}; pseudo-rehearsal \citep{robins1995pseudo};
generative replay \citep{shin2017generative}; sleep-like replay in networks
\citep{tadros2022sleep}. SEAL \citep{zweiger2025seal} has a model write its
own finetuning data. Self-consistency \citep{wang2023selfconsistency} is our
voting rule.

\section{Limitations}
One seed per arm; several effects sit within one to two standard errors.
Corpora of 60 facts make each corpus-level percentage a handful of facts. The
answer vocabulary has ten names. The facts are synthetic and the paraphrase
probes are light. The fast tier is a bounded cache, so facts older than one
window rely on sleep. The XL model is worse at ordinary language modeling than
the matrix-memory model (3.51 against 3.34), and a cache bias did not close the
gap. Cue tuning trains the exact cue format the night uses; held-out subjects
and per-sample answers rule out memorized facts, but not that the read is
specific to this format. The best loop keeps exactly the pre-registered 50\%,
from one seed. Scale is 134M parameters and 393M tokens per round;
paraphrase-invariant retrieval at this size is capacity-limited (30\% in-context
ceiling on our hardest probe level). Lifetime tables are trained by
gradient descent; one-shot writes into the fast-weight memory across documents
failed (0\%).

\section{Conclusion}
A small model can learn by day and consolidate by night with no ground truth
after the day itself: it dreams about what it read, votes on its dreams, sleeps
on them, and keeps the result. The stage that failed first was not storage
but the step from memory to a dream: the cache held the facts and the model
could not answer a cue from it until it was trained to. After that it learns
more than half of a day and keeps half of that. The weak stage now is
rehearsal, where old facts are re-dreamed from imperfect weights. The two lessons we would pass on are methodological. Measure the
backbone's own ceiling on every probe before setting a bar. And when a new
memory module stores information but does not use it, check whether the model
already has a retrieval circuit you could hand the memory to.

\bibliographystyle{plainnat}
\bibliography{references}
\end{document}
